\chapter{如何选课}

如何选课？“国内绝大部分大学的本科教学，不是濒临崩溃，而是早已崩溃。”——《上海交通大学生存手册》这本册子即使已流传近二十年仍在各方面有可取之处，可供一读。

与中学里有明确的课程表、部分高校各班集体分配必修课不同，除了某些专业课因人数较少只开一个班外，我校的课程都是可以进行自由选择的。这意味着每名同学都能拥有与他人不同的课表、在不同老师的课堂学习同一门课。明智的选课可能是一个学期幸福的开始，接下来向各位新同学简单介绍我校选课的方法与技术要点。

\section{选课前认真阅读培养手册}

培养手册是每名同学选课的重要指导，其中会明确指出毕业所需的各类课程学分、各专业的修课要求。

我校每一年级的培养手册都可能会有一定变化，因此请在拿到手册后、选课前认真阅读培养手册。带着以下问题阅读有助于你更好地理解：

(1)如果我要顺利毕业，我要修读哪几类课程？它们的学分要求分别是怎样的？（请注意，学分不是“分”，它可以视为课程的重要性或者强度指标。所以请不要问“驾照能不能加学分”这样的问题了。）

(2)我想转到（或留在本专业）我校的哪个/哪些专业？这些专业对我的公共基础课选课有要求吗？

(3)每一类课程的具体构成大概是怎样的？我这个学期必须要上哪些课？请牢记，认真阅读培养手册是顺利完成毕业学分清查、（如果有意愿）顺利次辅修的关键所在。如果你对选课有任何困惑，请一定要阅读培养手册。（值得指出，培养方案虽对选课很重要，但切勿盲从培养方案、只盯着学校划定的课程去学习。）

\section{结合选课单和选课指南认真研究}

在我校 sep 综合信息平台

的“选课系统”中，你可以找到每个学期的开课单和各个年级的选课指南（建议认真熟悉选课系统）。开课单上的信息包括：老师、时间、地点、考核方式等；而选课指南是“说明书”，会告诉你选课各阶段的任务和时间、部分课程安排可能需要注意的微妙之处。

如果想排出一份令你满意的课表，那么在研究选课文件时，以下问题是需要考虑的：

(1)开课老师。在我校，即使几门课程的名称完全一致，只要教学老师不同，其教学方式、考核方式、给分情况都可能是完全不一样的。

以线性代数为代表的课程会出现很明显的“一数各表”的情况。在选课之前提前做好功课，向学长学姐询问老师情况，有助于你选到适合你的课程。

在这方面，每到选课临近的一两周，我校“果壳玉泉跳蚤市场”群就会高强度讨论各门课程老师教学情况。如果愿意了解，可以多多提问、经常浏览、翻阅聊天记录、用老师的名字/姓名首字母（我校特有的指代某人的方法）查找相关情况。

但请注意，我要强调国科大的一大原则“小马过河”。每位同学对同一个课程（也有可能是同一种食物、同一种商品……）都会有不一样的评价，这也很正常，毕竟每个人的性格和需求不同。但是，请不要因为某位前辈给你带来的错误印象而心生埋怨，因为有可能 ta真是这么觉得的。

真正的情况还需要各位新同学亲身实践、眼见为实。包括我所给出的所有建议，都需要向新同学指出：“小马过河”！

(2)课程时间和地点。选课系统会禁止你选择时间冲突的两门课程，因此请在选课之前简单把你的选课方案打一个草稿。

可以准备几套方案以备不测，确保你选课舒心满意。此外，公共必修课会安排专门的习题课，请也要将习题课时间考虑在内。

选课时可以考虑课程时间与自己的作息习惯。如果你喜欢熬夜或者赖床，那么你可能要考虑早八的个数。同时，如果你想要有供你灵活调动的大块时间，或者你未来可能需要拜访中关村/奥运村区域研究所的导师，那么我建议尝试空出一整个工作日半天。

如果你是放假休息的忠实爱好者，你还可以研究学期的调休情况，力争在调休日少上几节课，就约等于多休息一天。

马院的思政类课程、外院的外语类课程有时认为放假期间的课也得补课，所以如果你反感补课，可以尝试着操作一下。当然，我不知道是否存在以上操作。

地点上，大一同学的课程基本集中在阶梯教室区域。只需要确保你能在课间休息时走到目标教室即可。在玉泉路校区，这基本不难。

(3)量力而行。需要指出，我校有些课程容量较大，可能上了一节课需要两节课甚至更多的时间去消化，因此请注意不要把自己的课表塞得太满。虽然我校准备了开学调课、中期退课等一系列机制为您留下了后路，但盲目选课也可能会让其他同学选不上想要的课程，建议您三思。例如：我校对每学期选课学分上限有硬性规定；外语选修课如果中期退课，那么之后您将失去外语选修课的优先选课权。

对于数学课程的 AB 班、物理课程的 4 学分班和 3 学分班的选择、物理先修课程，也建议同学们量力而行。可以试着读一读课程用书，(如果没买也可以网上找，我校一些 QQ 群和 Z-Library 之类的网站可能会有电子版)或者向学长学姐咨询。值得指出的是，某些数学 B 班的难度可能也较高，有前辈称之为“Super B(SB)班”，这方面可以在选课周前咨询学长学姐。但是暑假就别急了。

(4)考虑课程的考核方式，有助于你有针对性地学习。例如，我校思政类课程不需要考试，考核以平时分与论文构成。了解这些有助于你日后的学习。对于不同种类课程，合适的配比调剂也可能会让你的学习更舒适。

\section{抢课}

如果想要尽量实现你的课表规划，请一定要牢记抢课时间，调整好网络，在指定时间进入选课系统开抢即可。勾上你想要的课，输入口算题验证码即可实现抢课。如果你需要抢的课比较多，那么分批提交、先抢最抢手的，也是一种方法。

特别地，大一上学期的数学、物理公共课实施“基于选课的分配制度”，如果选择该班的同学超额，那么将会抽签分配同学到同一层次的其他班级。大学英语、大学写作课程也由学校分配。再次提醒，届时请您认真阅读培养手册、选课指南等文件并认真执行。

\section{如果遇到困难，及时调整}

选课过程一环扣一环，没有选到

自己想要的课实在正常不过，所以可以多准备几套备选方案。同时，你可以选择同学互相调课互通有无（可以选个夜深人静的时刻，同时退课互相选对方空出来的名额，以免被人截胡）。如需增选/调课，也可以向本科部负责老师提出。

具体流程可以在跳蚤市场群向前辈们询问。选课后，也请关注选课系统可能公布的课时变更/教材变更。

如果某一课程选课后让你实在无法接受、想要退课，可考虑：(1)学期前几周自由退课重选；(2)期中前的中期退课；(3)数学类课程在大一上学期期末允许 A 班换至 B 班等方法。如有任何问题，也可向辅导员老师询问，只要措辞得体，老师们会尽力帮助你解决问题。

\section{批判看待“谁给分好就选谁”的思维}

给分好不好本身就是

主观性极强的话题，且给分情况与教学质量不构成必然关系。此处我引用 21 级吉骏雄学长的话（之后可能会在绩点篇再重复一遍）：

这几天经常看到新生说一些“不卷 GPA 就没有保研，不保研就没有研究生读，没研究生读就没有好工作，没有好工作就没有好未来”的话。我想说的是，虽然 GPA 有用，但它是你人生路上的一个副产物。

无论你是用什么方法，譬如认真掌握了专业知识，譬如考前摸摸课本，譬如找人要去年期末考试试卷，发现题目全部都一样，这都只是副产品。如果你的个人能力够，这个 GPA 其实并不那么重要；就像今年寄系的推免名额都没用完，虽然原因不完全是你想象的那样罢了。

在这种学校，不想着怎样利用国科大的资源，怎样提高自己的学术水平，或者不想着怎么玩遍中国，怎么以最低的代价速通本科阶段，或者不想着自己怎样天天过得开心、过得痛快，而是张嘴 GPA 闭嘴内卷，每天除了做没有意义的题目就是互相以卖菜的形式攀比，那我只能说你的人生可能失去了最精彩的几年。

\section{两个擅自一个退路}

22 级黄铮舸学长提出过“两个擅自一个退路”，供各位新同学参考：不要擅自期待课程内容与考试成绩，不要擅自对课程及其成绩破防。选课时要给自己留好退路。

以上是本人所能想到的“选课”有关的所有注意事项。祝每位新同学顺利选上自己心仪的课程！
